% ПРИЛОЖЕНИЯ, без номер.
% По А.3: изходен текст на програми, формати на структури от данни, екранни
% форми, получени резултати. Използват се \subsection* без номер, за да не се
% продължи номерацията на параграфите от Раздел VI.
\section*{ПРИЛОЖЕНИЯ}
\label{sec:appendix}
\addcontentsline{toc}{section}{ПРИЛОЖЕНИЯ}

\subsection*{Приложение А. Формат на файловете с индекси}

\TODO{двоичният формат на един индексен файл: заглавен блок, ред на записите,
подравняване и признак за версия}

\subsection*{Приложение Б. Формат на речника на термините}

\TODO{представяне на двете посоки на съпоставянето и начин на кодиране на вида на
термина в идентификатора}

\subsection*{Приложение В. Граматика на поддържаното подмножество на SPARQL}

\TODO{изброяване на продукциите от граматиката в спецификацията, които се
разпознават, и на онези, които се отхвърлят с изрична грешка}

\subsection*{Приложение Г. Изходен текст на избрани модули}

\TODO{листинги на анализатора на заявки и на оператора за съединяване чрез сливане,
след като бъдат написани}

\subsection*{Приложение Д. Заявки, използвани при измерванията}

\TODO{пълният текст на всяка заявка от работното натоварване по Раздел V}

\subsection*{Приложение Е. Указания за изграждане и стартиране}

\TODO{командите за изграждане и за пускане на тестовете, сверени с README на
хранилището с изходния текст}
